\section{Semantics}
\label{sec:smt-semantics}

As for the source language, the semantics of SMT developed in this work is
\emph{denotational}: each term is assigned a value in the set-theoretic domain of
the ZFLean framework of \Cref{ch:zflean}.
Concrete terms are interpreted after abstraction to PHOAS, as described in
\Cref{subsec:smt-syntax-change-repr}.

\subsection{Interpretation of SMT types}
\label{subsec:smt-type-interp}

We first interpret SMT types as ZFC sets, denoted \(\interpZ{\cdot}\).

\begin{definition}[Type interpretation][SMT.SMTType.toZFSet][SMT/Semantics.html\#SMT.SMTType.toZFSet]
  \label{def:smt-type-interp}
  The function \(\interpZ{\cdot} : \mathsf{SMTType} \to \mathsf{ZFSet}\) is
  defined by structural recursion:
  \[
    \begin{array}{rclcrcl}
      \interpZ{\boolS}              & = & \Zbool, & \qquad &
      \interpZ{\intS}               & = & \Zint,                                  \\[2pt]
      \interpZ{\unitS}              & = & \{\emptyset\}, & \qquad &
      \interpZ{\alpha \timesS \beta}& = & \interpZ{\alpha} \timesZ \interpZ{\beta}, \\[2pt]
      \interpZ{\alpha \toS \beta}   & = & {\interpZ{\beta}}^{\interpZ{\alpha}}, & \qquad &
      \interpZ{\optS\tau}           & = & \regnotation{\optZ}(\interpZ{\tau}),
    \end{array}
  \]
  where \({\interpZ{\beta}}^{\interpZ{\alpha}}\) denotes the set of total
  functions (functional graphs) from \(\interpZ{\alpha}\) to \(\interpZ{\beta}\),
  and \(\optZ\) the encoding of options in ZFC (see \Cref{def:zflean-sum}).
\end{definition}

As on the source side, the set-theoretic interpretation is appropriate only if
SMT types are inhabited.
SMT types admit canonical inhabitants at the abstract term level, playing a dual role that
the carrier embedding \(\lift{\cdot}\) of \Cref{def:btype-to-term} plays on the
source side.

\begin{definition}[Default inhabitants][SMT.SMTType.default][SMT/Semantics.html\#SMT.SMTType.default]
  \label{def:smt-default}
  The default inhabitant \(\mathsf{default}(\tau)\) of an SMT type \(\tau\) is defined by
  structural recursion:
  \[
    \mathsf{default}(\tau) \triangleq
    \begin{cases}
      0 &\text{ if } \tau = \intS,\\
      \mathsf{false} &\text{ if } \tau = \boolS,\\
      \aop{\unitvalS} &\text{ if } \tau = \unitS, \\
      \aop{\noneS}_{\alpha} &\text{ if } \tau = \optS\alpha,\\
      \pairSa{\mathsf{default}(\alpha)}{\mathsf{default}(\beta)} &\text{ if } \tau = \alpha \timesS \beta,\\
      \aop{\lamS}\, (\_ : \alpha) \cdot \mathsf{default}(\beta) &\text{ if } \tau = \alpha \toS \beta.
    \end{cases}
  \]
\end{definition}

\begin{lemma}[Nonemptiness][SMT.SMTType.toZFSet\_nonempty][SMT/Semantics.html\#SMT.SMTType.toZFSet_nonempty]
  \label{thm:smt-type-nonempty}
  Interpretations of SMT types are non-empty: \(\forall \tau,\, \interpZ{\tau} \neq \emptyset\).
\end{lemma}
\begin{proof}
  By structural induction on \(\tau\), and \Cref{def:smt-default} provides a witness in each case.
\end{proof}

\subsection{Semantic domain}
\label{subsec:smt-sem-domain}

We now define the semantic domain of SMT terms. As for the source language
(\Cref{subsec:b-sem-domain}), we constrain the domain to carry intrinsic SMT types,
so that denotations are defined only for \emph{well-typed} values. This
correct-by-construction design removes all type-correctness side conditions from
the denotation clauses, leaving only the genuine partiality of partial
operations. The rationale for this type-carrying domain and the correct-by-construction
discipline it follows---together with the discarded alternative of denoting typing
derivations instead---is discussed in \Cref{subsec:b-denotation-design} and applies
unchanged here.

\begin{definition}[Semantic domain][SMT.Dom][SMT/Semantics.html\#SMT.Dom]
  \label{def:smt-dom}
      The semantic domain is the dependent sum
  \[
    \SDom \triangleq
    \sum_{X : \mathsf{ZFSet}}\sum_{\tau : \mathsf{SMTType}}\, X \in \interpZ{\tau}.
  \]
\end{definition}

Elements of \(\SDom\) are dependent triples, denoted \(\denval{X}[\tau]\),
consisting of a ZFC set \(X\), an SMT type \(\tau\), and a proof
\(\pf{\tau}{X} : X \in \interpZ{\tau}\). As before, we write ``\(\irrelprf\)'' for
the proof component when it is irrelevant.

\subsection{Denotation}
\label{subsec:smt-denote}

The denotation is a partial function over PHOAS terms parameterized by the
semantic domain, following the idea of \Cref{subsec:phoas-for-semantics}.

\begin{definition}[Denotation][SMT.denote][SMT/Semantics.html\#SMT.denote]
  \label{def:smt-denote}
    The denotation function is defined by structural recursion over PHOAS terms as
  the following partial function
  \[
    \regnotation{\interpS{\cdot}} : \mathsf{PHOAS.Term}\,\SDom \rightharpoonup \SDom.
  \]
\end{definition}

The defining equations are detailed below, written \(\interpS{t} \triangleq
\denval{X}[\tau]\); in each case the membership proof \(\pf{\tau}{X}\) is
constructed alongside the value.

\paragraph{Variables and constants.}
Variables denote themselves: for any variable \(\var{v} : \SDom\),
\(\interpS{\var{v}} \triangleq \var{v}\). For an integer \(n\), a Boolean \(b\), and the unit
value,
\[
  \interpS{n} \triangleq \denval{\zop{n}}[\intS],\quad
  \interpS{b} \triangleq \denval{\zop{b}}[\boolS],\quad
  \interpS{\aop{\unitvalS}} \triangleq \denval{\emptyset}[\unitS],
\]
where \(\zop{n}\) and \(\zop{b}\) are the ZFC encodings of integers and Booleans.
\begin{proofconstruct}
  \(\pf{\intS}{\zop{n}} : \zop{n} \in \Zint\) and \(\pf{\boolS}{\zop{b}} : \zop{b} \in \Zbool\)
  are supplied by {\small\leancodeptr{ZFSet.mem\_ofInt\_Int}{ZFLean/Integers.html\#ZFSet.mem_ofInt_Int}}
  and {\small\leancodeptr{ZFSet.ZFBool.mem\_ofBool\_𝔹}{ZFLean/Booleans.html\#ZFSet.ZFBool.mem_ofBool_𝔹}}
  respectively; \(\pf{\unitS}{\emptyset} : \emptyset \in \{\emptyset\}\) trivially holds.
\end{proofconstruct}

\paragraph{Pairs and projections.}
Given \(\interpS{t} = \denval{X}[\alpha]\) and \(\interpS{s} = \denval{Y}[\beta]\),
\[
  \interpS{\pairSa{t}{s}} \triangleq \denval{\pairZ{X}{Y}}[\alpha \timesS \beta],
\]
using Kuratowski pairs (see \Cref{subsec:zflean-sets}). Projections are defined when the argument has a pair
type: given \(\interpS{t} = \denval{T}[\alpha \timesS \beta]\),
\[
  \interpS{\fstSa t} \triangleq \denval{\fstZ(T)}[\alpha], \qquad
  \interpS{\sndSa t} \triangleq \denval{\sndZ(T)}[\beta],
\]
where \(\fstZ,\, \sndZ\) are the set-theoretic projections of \Cref{def:zflean-projections}.
\begin{proofconstruct}
  \(\pf{\alpha \timesS \beta}{\pairZ{X}{Y}} : \pairZ{X}{Y} \in \interpZ{\alpha} \timesZ \interpZ{\beta}\)
  is given by {\small\leancodeptr{ZFSet.pair\_mem\_prod}{Mathlib/SetTheory/ZFC/Basic.html\#ZFSet.pair_mem_prod}}
  (see \Cref{lst:pair-proj}) composed with \(\pf{\alpha}{X}\) and \(\pf{\beta}{Y}\).
  For projections, \(\pf{\alpha}{\fstZ(T)}\) and \(\pf{\beta}{\sndZ(T)}\) are extracted from
  \(\pf{\alpha \timesS \beta}{T}\) using
  {\small\leancodeptr{ZFSet.π₁\_pair}{Mathlib/SetTheory/ZFC/Basic.html\#ZFSet.π₁_pair}} and
  {\small\leancodeptr{ZFSet.π₂\_pair}{Mathlib/SetTheory/ZFC/Basic.html\#ZFSet.π₂_pair}} (see \Cref{lst:pair-proj}).
\end{proofconstruct}

\paragraph{Arithmetic and comparison.}
Given \(\interpS{t} = \denval{X}[\intS]\) and \(\interpS{s} = \denval{Y}[\intS]\)
(otherwise the denotation is undefined), arithmetic lifts the corresponding ZFC
operations:
\[
  \interpS{t \mathbin{\aop{\addS}} s} \triangleq \denval{X \addZ Y}[\intS],\quad
  \interpS{t \mathbin{\aop{\subS}} s} \triangleq \denval{X \subZ Y}[\intS],\quad
  \interpS{t \mathbin{\aop{\mulS}} s} \triangleq \denval{X \mulZ Y}[\intS],
\]

\begin{proofconstruct}
  All membership proofs are instances of
  {\small\leancodeptr{overloadBinOp\_mem}{Extra/ZFC_Extra.html\#overloadBinOp_mem}},
  which provides closure of overloaded binary operations: \(\Zint\) is closed under
  \(\addZ\), \(\subZ\), and \(\mulZ\).
\end{proofconstruct}

\paragraph{Boolean connectives.}
Given \(\interpS{t} = \denval{X}[\tau]\) and \(\interpS{s} = \denval{Y}[\sigma]\),
the primitive connectives lift their ZFC counterparts in the two-element Boolean algebra \(\Zbool\):
\begin{align*}
  \interpS{t \mathbin{\aop{\eqS}} s}  & \triangleq \denval{X = Y}[\boolS] && \text{when } \tau = \sigma, \\
  \interpS{t \mathbin{\aop{\andS}} s} & \triangleq \denval{X \andZ Y}[\boolS] && \text{when } \tau = \sigma = \boolS, \\
  \interpS{\aop{\notS} t}   & \triangleq \denval{\negZ X}[\boolS] && \text{when } \tau = \boolS,\\
  \interpS{t \mathbin{\aop{\leS}} s} &\triangleq \denval{X \leqZ Y}[\boolS] && \text{when } \tau = \sigma = \intS.
\end{align*}
Disjunction, implication, and the existential quantifier are interpreted through
their definitions as derived forms (\Cref{def:smt-derived}).
\begin{proofconstruct}
  All membership proofs follow from
  {\small\leancodeptr{overloadBinOp\_mem}{Extra/ZFC_Extra.html\#overloadBinOp_mem}} for
  binary operations and
  {\small\leancodeptr{overloadUnaryOp\_mem}{Extra/ZFC_Extra.html\#overloadUnaryOp_mem}}
  for negation, each asserting that the respective ZFC operations are closed on \(\Zbool\),
  and \(\leqZ\) maps a pair of integers to \(\Zbool\).
\end{proofconstruct}

\paragraph{Options and unit.}
Given \(\interpS{t} = \denval{T}[\tau]\), options are interpreted by the tagged encoding
of \Cref{def:smt-type-interp}:
\[
  \interpS{\someSa t} \triangleq \denval{\someZ\,T}[\optS\tau],\qquad
  \interpS{\aop{\noneS}_{\tau}} \triangleq \denval{\noneZ}[\optS\tau].
\]
Conversely, given \(\interpS{t} = \denval{T}[\optS\tau]\),
\[
  \interpS{\theSa t} \triangleq \denval{\theZ\,T}[\tau].
\]
Here \(\theZ\) is the selector of the option encoding (\Cref{def:zflean-sum}): it satisfies
\(\theZ(\someZ X) = X\), and on \(\noneZ\) it returns a designated element of the carrier
\(\interpZ{\tau}\)---well-defined because \(\interpZ{\tau}\) is non-empty
by \Cref{thm:smt-type-nonempty}. It is therefore \emph{total}, faithfully reflecting
SMT-LIB's treatment of a datatype selector applied to a non-matching constructor as
underspecified rather than undefined.
For the unit type \(\unitS\),
\[
  \interpS{\aop{\unitvalS}} \triangleq \denval{\emptyset}[\unitS].
\]

\begin{proofconstruct}
  The proofs are straightforward and follow from the definitions of \(\someZ\), \(\noneZ\),
  \(\theZ\), and \(\interpZ{\unitS} \triangleq \{\emptyset\}\).
\end{proofconstruct}

\paragraph{Distinctness.}
Distinctness of a family is the iterated, pairwise conjunction asserting that no
element equals any later one, folded in \(\Zbool\), abbreviated here as the following
function \(\mathbb{d}\):
\begin{equation*}
  \text{for any family of ZFC sets } \vec z = (z_1, \cdots, z_n),\qquad \mathbb{d}(\vec z) \coloneq \bigandZ_{1 \leq i < j \leq n} z_i \neq z_j.
\end{equation*}
Let \(\vec t = (t_1, \dots, t_n)\) be a family of SMT terms such that every \(t_i\) has a denotation.
The denotation of the distinctness predicate is then defined as
\[
  \interpS{\distinctSa \vec t} \triangleq \denval{\mathbb{d}((\interpS{t_i})_i)}[\boolS].
\]
\begin{proofconstruct}
  The proof that \(\mathbb{d}((\interpS{t_i})_i) \in \Zbool\) is established by structural
  induction on the family \(\vec t\), using that \(\andZ\) is a closed operation on \(\Zbool\);
  the base case \(\topZ \in \Zbool\) is immediate.
\end{proofconstruct}

\paragraph{Application.}
Application is the only genuinely partial operation, and is
interpreted by graph lookup. Given \(\interpS{f} = \denval{F}[\sigma \toS \tau]\) and
\(\interpS{x} = \denval{X}[\xi]\), the denotation is defined when the types
match (\(\sigma = \xi\)), \(F\) is a functional graph from
\(\interpZ{\sigma}\) to \(\interpZ{\tau}\), and \(X\) lies in its domain; letting
\(Y\) be the unique value with \(\pairZ{X}{Y} \in F\), the denotation is given by
\[
  \interpS{\appSa{f}{x}} \triangleq \denval{Y}[\tau].
\]
\begin{proofconstruct}
  Since by assumption, \(F \in \interpZ{\sigma \toS \tau} = {\interpZ{\tau}}^{\interpZ{\sigma}}\),
  \(F\) is a total function with domain \(\interpZ{\sigma}\).
  Such a witness \(Y\) therefore exists and is unique by definition, and the proof
  \(\pf{\tau}{Y} : Y \in \interpZ{\tau}\) naturally follows from the membership of
  \(\pairZ{X}{Y}\) in \(F\):
  \[
    \pairZ{X}{Y} \in F \in {\interpZ{\tau}}^{\interpZ{\sigma}} \subseteq \interpZ{\tau} \timesZ \interpZ{\sigma}.
    \vspace{-\baselineskip}
  \]
\end{proofconstruct}

\begin{remark}
  \label{rem:smt-wd}
  This is the only well-definedness side condition in the SMT semantics. The
  source language additionally requires finiteness for cardinality (\Cref{sec:b-wd});
  SMT-LIB has no such operators in the fragment
  considered, and its function spaces are total, so the denotation is partial
  only at application---and even there, the adequacy results below show that the
  side condition is automatically satisfied for well-typed terms whose function
  argument is supplied within its domain.
\end{remark}

\paragraph{Binders.}
Binder constructs rely on the PHOAS representation, exactly as in
\Cref{subsec:b-denote}. For a binder \(\aop{\lamS}\vec\tau \cdot e\) of arity
\(n\), let \(e\) denote its PHOAS scope and write
\(\mathcal{D} \coloneq \interpZ{\tau^{\times}}\) for the carrier of the argument
tuple. Each \(z \in \mathcal{D}\) decomposes into
components \(z_i \in \interpZ{\tau_i}\) and induces a family of domain values
\(\vec z = (\denval{z_i}[\tau_i])_i\); the scoped denotation of the body is the
function on \(\mathcal{D}\) obtained by interpreting the body at \(\vec z\).

\emph{Lambda abstraction} is denoted by the graph of the scoped body. We write
\(\Phi(z)\) for the (set component of the) body denotation at the tuple \(z\), and
\(\gamma\) for its type, \ie
\[
  \interpS{e(\vec z)} = \denval{\Phi(z)}[\gamma].
\]
Notice that \(\gamma\) depends on the argument \(z\), and such a definition does not provide a
return type for the lambda abstraction. We break this dependency using \Cref{def:smt-default} by
evaluating the body at the canonical default tuple \(\vec d \coloneq (\denval{\mathsf{default}(\tau_i)}[\tau_i])_i\):
\[
  \interpS{e(\vec d)} = \denval{\Phi(\vec d)}[\gamma][\irrelprf].
\]
Evaluating at this fixed tuple makes the choice of \(\gamma\) homogeneous across all arguments
(\ie \(\gamma\) is the same for all \(z \in \mathcal{D}\)), and no longer dependent on the argument \(z\).

\begin{note}
  Recall that in this setting---contrary to the source language---functions are total, so the body
  denotation is defined for every argument in the domain. In particular, \(\mathcal{D}\) is
  non-empty by \Cref{thm:smt-type-nonempty}, so the default tuple \(\vec d\) is well-defined.
\end{note}

The denotation of the lambda abstraction is then given by the following set-level abstraction
(see \Cref{subsec:zflean-func-eval-lambda}), where
\(
  \Lambda \coloneq \lambdaZ{\mathcal{D}}{\interpZ{\gamma}}{z}{\Phi(z)}
\):

\[
  \interpS{\aop{\lamS} \vec\tau \cdot e}
  \triangleq
  \denval{\Lambda}[\tau^{\times} \toS \gamma][\pf{}{\lamS}].
\]

\emph{Universal quantification} is denoted by the bounded infimum of the body values
over the domain, taken in the two-element Boolean algebra:
\[
  \interpS{\aop{\forallS} \vec\tau \cdot P}
  \triangleq
  \denval{\bigcap \{\, y \in \Zbool \mid \exists z \in \mathcal{D},\ y = \Phi(z) \,\}}[\boolS][\pf{}{\forallS}],
\]
keeping the notation \(\Phi(z)\) for the body denotation at the tuple \(z\).

\begin{note}
Again, the infimum is well-defined since \(\mathcal{D}\) is non-empty, which ensures that the set of body values
is likewise non-empty. Consequently, \(\forallS\) requires no special treatment for the empty domain,
in contrast to the source language.
\end{note}

\emph{Existential quantification} is denoted by the derived form
\(\aop{\existsS} \vec\tau \cdot P \triangleq \notS(\forallS \vec\tau \cdot \notS P)\).

\begin{proofconstruct}
  For lambda abstraction, \(\pf{}{\lamS} : \Lambda \in {\interpZ{\gamma}}^{\mathcal{D}}\)
  is supplied by the fact that
  \(\Phi(z) \in \interpZ{\gamma}\) for every \(z \in \mathcal{D}\), and the fact that set-level
  abstractions are total functions (\Cref{thm:lambda}; in Lean, theorem
  {\small\leancodeptr{ZFSet.mem\_funs\_of\_lambda}{ZFLean/Functions.html\#ZFSet.mem_funs_of_lambda}}).
  For universal quantification, \(\pf{}{\forallS}\) follows from \(\Zbool\) being closed under
  the infimum (or big intersection), and the fact that \(\Phi(z) \in \Zbool\) for every
  \(z \in \mathcal{D}\)---given by theorem 
  {\small\leancodeptr{ZFSet.sInter\_sep\_subset\_of\_𝔹\_mem\_𝔹}{ZFLean/Booleans.html\#ZFSet.sInter_sep_subset_of_𝔹_mem_𝔹}}.
\end{proofconstruct}

\subsection{Basic properties}
\label{subsec:smt-sem-properties}

The denotation enjoys the same two fundamental properties as on the source side:
\emph{type adequacy} and \emph{totality} for well-typed terms, together
resembling the soundness and completeness of the semantics with respect to the
type system.

\paragraph{Type adequacy.}
Adequacy states that whenever a term denotes, the intrinsic SMT type carried by its
denotation coincides with the type assigned by any typing derivation. It rests on
the determinism of the type system (\Cref{thm:smt-typing-det}).

\begin{theorem}[Denotation typing-soundness][SMT.PHOAS.denote\_welltyped\_eq][SMT/Reasoning/Lemmas.html\#SMT.PHOAS.denote_welltyped_eq]
  \label{thm:smt-denote-type-det}
  Let \(t : \mathsf{PHOAS.Term}\,\SDom\) be a term admitting a typing
  \(\satype{\Gamma}{t}{\tau}\) under a well-formed context \(\Gamma\). If
  \(\interpS{t} = \denval{T}[\sigma][\irrelprf]\), then \(\tau = \sigma\).
\end{theorem}
\begin{proof}
  By induction on \(t\). The key observation is that the SMT type carried by a
  denotation is constructed from the subterm denotations by the very same
  combinator that the typing rules use to combine subterm types. The induction
  therefore reduces to inverting each constructor and identifying the carried
  type with the assigned one.
\end{proof}

This result calls for two comments.
First, it is the semantic counterpart of typing determinism: just as a well-typed
term has a unique static type, its denotation---whenever defined---carries exactly
that type. Equivalently, a term cannot be denoted by a value whose intrinsic type
differs from the one assigned by its \(\saop{\vdash}\)-derivation. The guarantee is
substantive in a set-theoretic setting, where a single set may inhabit several
types; the theorem ensures that the static type alone fixes the type carried by the
denotation, so that \(\interpS{\cdot}\) never ascribes conflicting types to a term.

Second, the result implies that, over well-formed contexts, the static type of a
term is independent of the context in which it is derived. This may seem
surprising: one might expect the same variable \(\var{x}\) to be typed differently under
different contexts, as in
\[
  \satype{\{ \var{x} \mapsto \intS\}}{\var{x}}{\intS} \quad\text{and}\quad
  \satype{\{ \var{x} \mapsto \boolS\}}{\var{x}}{\boolS},
\]
leaving the type carried by \(\interpS{\var{x}}\) ambiguous. The ambiguity is
only apparent. By definition \(\interpS{\var{x}} \triangleq \var{x}\), where \(\var{x}\)
is a triple \(\denval{X}[\tau][\irrelprf]\) carrying a fixed intrinsic type
\(\tau\). Well-formedness (\Cref{def:smt-wftc}) requires every context to assign
\(\var{x}\) precisely this intrinsic type, so the two contexts above cannot both be
well-formed---one would demand \(\tau = \intS\), the other \(\tau = \boolS\)---and
the denotation unambiguously carries \(\tau\). More generally, the variable rule is
the only context-sensitive rule, and well-formedness pins every variable to its
intrinsic type; the static type of a well-typed PHOAS term is therefore determined
by the term alone.


\paragraph{Totality.}
Totality is the converse existence statement: well-typed terms have a denotation,
of the expected type. As on the source side, the abstraction of a concrete term
has a denotation provided the typing context is well-formed and the renaming
covers the term's free variables.

\begin{theorem}[Denotation typing-completeness][SMT.RenamingContext.denote\_exists\_of\_typing][SMT/Reasoning/Basic/DenotationTotality.html\#SMT.RenamingContext.denote_exists_of_typing]
  \label{thm:smt-denote-complete}
  Let \(\Gamma\) be a concrete typing context, \(t\) a concrete SMT term with
  \(\stype{\Gamma}{t}{\tau}\), and \(\Delta\) a renaming that covers the free variables of \(t\). 
  Then \(\aterm{t}{\Delta}\) has a denotation carrying type \(\tau\), \ie
  \[
    \exists\, T \in \interpZ{\tau},\; \interpS{\aterm{t}{\Delta}} = \denval{T}[\tau].
  \]
\end{theorem}
\begin{proof}
  By induction on the typing derivation \(\stype{\Gamma}{t}{\tau}\). We proceed by a case analysis divided into four groups.

  \paragraph{Variables and constants.}
  For variables, type-compatibility supplies a domain value \(\Delta(\var{v})\) of
  type \(\tau\) and abstraction maps \(\var{v}\) to it.
  Constants are immediate, with the membership
  proofs constructed in \Cref{subsec:smt-denote}.

  \paragraph{Compound non-partial constructors.}
  The induction hypothesis applied to each subterm yields a witness denotation
  whose type matches the corresponding typing premise.
  The witness set is then assembled by the relevant ZFC closure lemma.

  \paragraph{Application.}
  The induction hypothesis gives
  \(\interpS{f} = \denval{F}[\sigma \toS \tau]\) and
  \(\interpS{x} = \denval{X}[\sigma]\). Since \(F \in \interpZ{\sigma \toS \tau} =
  {\interpZ{\tau}}^{\interpZ{\sigma}}\) is a functional graph with domain
  \(\interpZ{\sigma}\)---which contains \(X\)---the well-definedness side condition
  of the application clause is therefore satisfied, and the unique \(Y\) with
  \(\pairZ{X}{Y} \in F\) provides the witness, with
  \(\pf{\tau}{Y} : Y \in \interpZ{\tau}\) following from
  \(\pairZ{X}{Y} \in F \subseteq \interpZ{\sigma} \timesZ \interpZ{\tau}\).

  \paragraph{Binders.}
  We treat the three binder cases uniformly. Inversion of the binder typing rule
  supplies, for some arity \(n \geq 1\), a type vector \(\vec\tau\) and a scoped body
  \(e\) well-typed under \(\Gamma[\vec{\var{v}}:\vec\tau]\) for every family of variables
  \(\vec{\var{v}}\) carrying the types \(\vec\tau\), with type \(\boolS\) for universal
  quantification and some \(\gamma\) for lambda abstraction.
  The argument proceeds in four steps.
    
  \smallskip\noindent\emph{Step 1: domain denotation.} The carrier \(\mathcal{D} = \interpZ{\tau^{\times}}\)
    is the product of the component carriers; each tuple
    \(z \in \mathcal{D}\) decomposes into components
    \(z_i \in \interpZ{\tau_i}\).
    
  \smallskip\noindent\emph{Step 2: body denotation at typed tuples.} Each \(z\) induces a family
    \(\vec z = (\denval{z_i}[\tau_i])_i\) of the prescribed types.
    The renaming context \(\Delta\) is extended to
    \(\Delta' \coloneq \Delta[\vec{\var{v}} \leftarrow \vec z]\);
    the induction hypothesis, applied to the body under \(\Delta'\)
    yields a body denotation of the right type, namely \(\boolS\)
    for universal quantification and \(\gamma\) for lambda abstraction.
    
  \smallskip\noindent\emph{Step 3: discharging the side conditions of the denotation clause.}
    For lambda abstraction, the denotation clause is guarded by two
    conditions: that the body has a denotation for every typed argument
    tuple, and that the body's type is uniform across all typed tuples.
    For universal and existential quantification only denotation definedness is needed,
    since the body always has type \(\boolS\) and uniformity is trivial.
    Step 2 discharges this condition in each case; type adequacy
    (\Cref{thm:smt-denote-type-det}) discharges the uniformity condition for
    lambda by identifying the body type uniquely with \(\gamma\), regardless
    of the tuple chosen.
    
  \smallskip\noindent\emph{Step 4: assembling the witness.} Writing \(\Phi\) for
    the scoped body denotation---\ie informally, for any family of sets \(\vec z\),
    \(\interpS{e((\denval{z_i}[\tau_i][\irrelprf])_i)} = \denval{\Phi(z)}[\gamma][\irrelprf]\)
    where \(\gamma\) matches the type of the body---the
    witness is assembled by separation: the graph
    \(\{\, \pairZ{z}{y} \in \mathcal{D} \timesZ \interpZ{\gamma} \mid y =
    \Phi(z) \,\}\) for lambda, which lies in
    \(\interpZ{\tau^{\times} \toS \gamma}\); the bounded infimum
    \(\bigcap\{\, y \in \Zbool \mid \exists z \in \mathcal{D},\ y = \Phi(z)
    \,\}\) for universal quantification, which lies in \(\Zbool\); and the negation
    of the corresponding universal denotation for the existential.
  
  In each case the constructed witness belongs to \(\interpZ{\tau}\) as expected,
  providing \(\pf{\tau}{T}\) and concluding the proof.
\end{proof}

This chapter has developed the formal target-language infrastructure for SMT-LIB 2.7.
\Cref{sec:smt-syntax} introduced the syntactic categories: a concrete, first-order
representation of terms that closely follows the surface syntax of SMT-LIB 2.7 and
constitutes the deep embedding emitted by BEer, and the PHOAS abstract representation used
for semantic reasoning.
\Cref{sec:smt-typing} gave typing rules for both representations, and established
several structural properties, including the determinism of the abstract type system
(\Cref{thm:smt-typing-det}), which underpins the adequacy result below.
\Cref{sec:smt-semantics} assigned denotations in the shared ZFC universe, proved type
adequacy (\Cref{thm:smt-denote-type-det}) and typing completeness
(\Cref{thm:smt-denote-complete}).

A key structural difference from the B chapter (\Cref{ch:b}) is that SMT function
spaces are total: \(\interpZ{\sigma \toS \tau} = {\interpZ{\tau}}^{\interpZ{\sigma}}\)
is the set of total functions, so the denotation is partial only at function application,
and even there only due to the type-matching guard---a guard that completeness shows is
automatically satisfied for well-typed terms. On the source side, by contrast, partiality
arises at every partial operator, \eg application and cardinality,
all requiring explicit well-definedness conditions.
